% STEMI 漏診案例分析教學講義
% 作者: 謝慕揚, MD, PhD, FESC
% 日期: 2026-01-06

\documentclass[12pt,a4paper]{article}

% Including essential packages for Chinese typesetting
\usepackage{xeCJK}
\usepackage{fontspec}
% Configuring Chinese font with Noto Serif CJK TC, fallback to SimSun
\IfFontExistsTF{Noto Serif CJK TC}{
  \setCJKmainfont{Noto Serif CJK TC}
  \setCJKsansfont{Noto Serif CJK TC}
  \setCJKmonofont{Noto Serif CJK TC}
}{\setCJKmainfont{SimSun}
  \setCJKsansfont{SimSun}
  \setCJKmonofont{SimSun}
}

% Including other necessary packages
\usepackage{geometry}
\usepackage{titlesec}
\usepackage{enumitem}
\usepackage{xcolor}
\usepackage{colortbl}
\usepackage{tcolorbox}
\usepackage{booktabs}
\usepackage{longtable}
\usepackage{array}
\usepackage{graphicx}
\usepackage{tikz}
\usepackage{fancyhdr}
\usepackage{multicol}
\usepackage{datetime2}
\usepackage{hyperref}
\usepackage{mdframed}
\usepackage{amsmath}
\usepackage{amssymb}
\usepackage{multirow}

% Defining colors
\definecolor{emergency_red}{RGB}{186,24,27}
\definecolor{emergency_blue}{RGB}{0,114,188}
\definecolor{emergency_gray}{RGB}{120,120,120}
\definecolor{highlight_yellow}{RGB}{255,250,205}
\definecolor{alert_orange}{RGB}{255,140,0}

\hypersetup{
    colorlinks=true,
    linkcolor=emergency_blue,
    citecolor=emergency_blue,
    urlcolor=emergency_blue,
    pdfborder={0 0 0}
}

% Defining page geometry
\geometry{
  top=2.5cm,
  bottom=2.5cm,
  left=2cm,
  right=2cm,
  headheight=25pt
}

% Adjusting topmargin to compensate for increased headheight
\addtolength{\topmargin}{-10pt}

% Setting title formats
\titleformat{\section}[block]{\Large\bfseries\color{emergency_red}}{\thesection}{10pt}{}
\titleformat{\subsection}[block]{\large\bfseries\color{emergency_blue}}{\thesubsection}{8pt}{}
\titleformat{\subsubsection}[block]{\normalsize\bfseries\color{emergency_gray}}{\thesubsubsection}{6pt}{}

% Defining custom environment for important notes
\newmdenv[
    linecolor=emergency_red,
    backgroundcolor=highlight_yellow,
    linewidth=2pt,
    topline=true,
    bottomline=true,
    leftline=true,
    rightline=true,
    innertopmargin=10pt,
    innerbottommargin=10pt,
    innerleftmargin=10pt,
    innerrightmargin=10pt
]{importantbox}

% Alert box for critical warnings
\newtcolorbox{alertbox}{
    colback=red!10!white,
    colframe=emergency_red,
    title=\textbf{嚴重警示},
    fonttitle=\bfseries\color{white},
    coltitle=white,
    colbacktitle=emergency_red,
    boxrule=1.5pt,
    arc=3pt,
    left=6pt,
    right=6pt,
    top=6pt,
    bottom=6pt
}

% Key point box
\newtcolorbox{keypoint}{
    colback=emergency_blue!5!white,
    colframe=emergency_blue,
    title=\textbf{重點提醒},
    fonttitle=\bfseries,
    boxrule=0.5pt,
    arc=3pt,
    left=6pt,
    right=6pt,
    top=6pt,
    bottom=6pt
}

% Learning point box
\newtcolorbox{learningbox}{
    colback=green!5!white,
    colframe=green!50!black,
    title=\textbf{學習要點},
    fonttitle=\bfseries,
    boxrule=0.5pt,
    arc=3pt,
    left=6pt,
    right=6pt,
    top=6pt,
    bottom=6pt
}

% Header and Footer
\pagestyle{fancy}
\fancyhf{}
\fancyhead[L]{\small\color{emergency_gray}急診 STEMI 漏診案例分析}
\fancyhead[R]{\small\color{emergency_gray}謝慕揚醫師}
\fancyfoot[C]{\thepage}
\renewcommand{\headrulewidth}{0.4pt}
\renewcommand{\footrulewidth}{0pt}

\begin{document}

% Title Page
\begin{titlepage}
\centering
\vspace*{2cm}

{\Huge\bfseries\color{emergency_red} 急診 STEMI 漏診案例分析}

\vspace{0.5cm}

{\LARGE\color{emergency_blue} 忽略心電圖電腦判讀警示的慘痛教訓}

\vspace{2cm}

{\Large 院內教學講義}

\vspace{1cm}

{\large 謝慕揚醫師 MD, PhD, FESC}

\vspace{0.5cm}

{\large 2026 年 1 月}

\vspace{2cm}

\begin{alertbox}
\centering
\textbf{本案例為真實醫療爭議案件}\\
已去識別化處理，僅供教學使用
\end{alertbox}

\vfill

{\small 新竹臺大分院 心臟血管中心}

\end{titlepage}

% Table of Contents
\tableofcontents
\newpage

%----------------------------------------------------------
\section{案例背景}
%----------------------------------------------------------

\subsection{病患基本資料}

\begin{longtable}{p{4cm}p{10cm}}
\toprule
\textbf{項目} & \textbf{內容} \\
\midrule
\endhead
年齡/性別 & 52 歲男性 \\
到院時間 & 2021/06/27 11:27 \\
到院方式 & 自行步入 (Walk-in) \\
主訴 (Chief complaint) & 胸痛 30 分鐘 \\
疼痛性質 & 悶痛 (Pressure-like chest pain) \\
疼痛評分 & 8/10 分 \\
持續時間 & 持續性疼痛 (Continuous) \\
過去病史 & 無特殊 (Nil) \\
檢傷分級 (Triage) & 第 2 級 \\
\bottomrule
\end{longtable}

\subsection{到院生命徵象}

\begin{longtable}{p{4cm}p{10cm}}
\toprule
\textbf{生命徵象} & \textbf{數值} \\
\midrule
\endhead
體溫 (BT) & 35.7°C \\
心跳 (PR) & 76/min \\
呼吸 (RR) & 20/min \\
血壓 (BP) & 137/99 mmHg \\
血氧飽和度 (SpO2) & 100\% \\
\bottomrule
\end{longtable}

\begin{keypoint}
病患呈現\textbf{典型 STEMI 臨床表現}：
\begin{itemize}[noitemsep]
\item 中年男性
\item 持續性胸悶痛 $>$ 30 分鐘
\item 疼痛評分高 (8/10)
\item 生命徵象穩定但有輕微高血壓
\end{itemize}
\end{keypoint}

\newpage
%----------------------------------------------------------
\section{時序事件分析}
%----------------------------------------------------------

\subsection{關鍵時間軸}

\begin{longtable}{p{2.5cm}p{11.5cm}}
\toprule
\textbf{時間} & \textbf{事件} \\
\midrule
\endhead
11:27 & 病患步入急診，主訴胸痛 30 分鐘，檢傷分級 2 級 \\
\midrule
\rowcolor{red!20}
\textbf{11:31:49} & \textbf{第一份 ECG：電腦判讀顯示}\newline
\colorbox{yellow}{\textbf{``Normal sinus rhythm, STE over inferior leads, suspected AMI''}} \\
\midrule
11:33 & 醫囑：抽血檢驗 (CBC, BMP, Troponin-T, D-dimer, Blood gas 等) \\
\midrule
11:34 & 醫囑：Complete EKG at 11:40 \\
\midrule
11:35 & 給予 Nitrostat Sublingual 0.6 mg \\
\midrule
11:38 & Chest X-ray (AP view) \\
\midrule
11:44:41 & 第二份 ECG：電腦判讀顯示 ``Sinus bradycardia, Abnormal ECG'' \\
\midrule
11:45 & 轉入觀察床位 A03 \\
\midrule
11:57 & 再給予 Nitrostat Sublingual 1\# \\
\midrule
12:22 & 醫囑：f/u TropT, EKG 14:30 (預約 2 小時後追蹤！) \\
\midrule
\rowcolor{red!20}
\textbf{14:13:24} & \textbf{第三份 ECG：電腦判讀再次顯示}\newline
\colorbox{yellow}{\textbf{``Sinus bradycardia, STE over inferior leads, suspected AMI''}} \\
\midrule
15:22:10 & 第四份 ECG：``資料品質不良...LVH with QRS widening'' \\
\bottomrule
\end{longtable}

\begin{alertbox}
\textbf{致命的疏忽：}電腦判讀\textbf{兩次}明確警示 ``\textbf{STE over inferior leads, suspected AMI}''，但醫療團隊\textbf{完全忽略}此警示訊息，未啟動心導管室！
\end{alertbox}

\newpage
\subsection{實驗室檢查結果}

\begin{longtable}{p{3cm}p{3cm}p{3cm}p{4cm}}
\toprule
\textbf{檢驗項目} & \textbf{11:39} & \textbf{14:17} & \textbf{正常值/判讀} \\
\midrule
\endhead
Troponin-T (ng/L) & 11.16 & \textcolor{red}{\textbf{56.43}} & $<$ 14 ng/L\newline\textcolor{red}{上升 5 倍！} \\
\midrule
tCO2 (mmol/L) & 23.8 & -- & 22-29 \\
\midrule
O2 Sat (\%) & 57.9 & -- & (靜脈血) \\
\midrule
D-dimer (mg/L) & 0.27 & -- & $<$ 0.5 \\
\midrule
Glucose (mg/dL) & 156 & -- & $<$ 100 (fasting)\newline壓力性高血糖 \\
\midrule
eGFR (mL/min) & 63.9 & -- & $>$ 90 \\
\bottomrule
\end{longtable}

\begin{keypoint}
\textbf{Troponin-T 動態變化：}
\begin{itemize}[noitemsep]
\item 初始值 11.16 ng/L（接近正常上限 14 ng/L）
\item 3 小時後追蹤值 56.43 ng/L（\textcolor{red}{\textbf{上升 5 倍}}）
\item 符合急性心肌梗塞的 Troponin 動態上升 ($>$ 20\% delta change)
\end{itemize}
\end{keypoint}

\newpage
%----------------------------------------------------------
\section{心電圖判讀分析}
%----------------------------------------------------------

\subsection{第一份 ECG (11:31:49) - 關鍵漏診時刻}

\begin{longtable}{p{4cm}p{10cm}}
\toprule
\textbf{項目} & \textbf{內容} \\
\midrule
\endhead
心率 (Heart rate) & 67 BPM \\
PR interval & 156 ms \\
QRS duration & 84 ms \\
電腦判讀 & \colorbox{yellow}{\textbf{Normal sinus rhythm}}\newline\colorbox{red!30}{\textbf{STE over inferior leads, suspected AMI}} \\
確認醫師 & OOO醫師 \\
\bottomrule
\end{longtable}

\subsection{Inferior STEMI 的 ECG 特徵}

\begin{longtable}{p{4cm}p{10cm}}
\toprule
\textbf{導程} & \textbf{預期變化} \\
\midrule
\endhead
\textbf{II, III, aVF} & \textcolor{red}{\textbf{ST elevation}} $\geq$ 1 mm\newline（本案例電腦判讀已偵測到） \\
\midrule
\textbf{I, aVL} & Reciprocal ST depression（對應性壓低） \\
\midrule
\textbf{V1-V3} & 可能有 ST depression（若合併 posterior wall involvement） \\
\midrule
\textbf{V4R} & 右心室梗塞時會有 ST elevation\newline\textcolor{emergency_blue}{（Inferior MI 應加做右側導程！）} \\
\bottomrule
\end{longtable}

\begin{alertbox}
\textbf{電腦判讀系統明確警示 ``suspected AMI''！}

這不是模糊的警示，而是\textbf{明確指出疑似急性心肌梗塞}。任何醫師看到此警示都應該：
\begin{enumerate}[noitemsep]
\item 立即親自判讀心電圖確認
\item 若有疑慮，立即會診心臟科
\item 考慮啟動心導管室
\end{enumerate}
\end{alertbox}

\newpage
%----------------------------------------------------------
\section{問題分析：為何會漏診？}
%----------------------------------------------------------

\subsection{系統性問題}

\begin{enumerate}[leftmargin=*]
\item \textbf{警示疲勞 (Alert fatigue)}
\begin{itemize}[noitemsep]
\item 電腦判讀系統經常有偽陽性
\item 醫師習慣性忽略電腦警示
\item 缺乏對高風險警示的優先處理機制
\end{itemize}

\item \textbf{認知偏誤 (Cognitive bias)}
\begin{itemize}[noitemsep]
\item 病患生命徵象穩定 → 低估嚴重度
\item 初始 Troponin 接近正常 → 排除 MI 的錯誤假設
\item 未考慮 Troponin 在 STEMI 早期可能正常
\end{itemize}

\item \textbf{流程問題}
\begin{itemize}[noitemsep]
\item ECG 確認醫師與主治醫師可能不同人
\item 缺乏強制性的警示回覆機制
\item 追蹤 ECG 排程過久（預約 2 小時後）
\end{itemize}
\end{enumerate}

\subsection{個人因素}

\begin{longtable}{p{5cm}p{9cm}}
\toprule
\textbf{問題} & \textbf{說明} \\
\midrule
\endhead
未仔細閱讀電腦判讀 & 只看心律，忽略 ``suspected AMI'' 警示 \\
\midrule
未親自判讀 ECG 波形 & 過度依賴電腦判讀或直接跳過 \\
\midrule
缺乏臨床警覺性 & 典型胸痛 + ECG 警示仍未聯想到 STEMI \\
\midrule
團隊溝通不足 & ECG 確認醫師與主治醫師間缺乏溝通 \\
\bottomrule
\end{longtable}

\newpage
%----------------------------------------------------------
\section{STEMI 診斷與處置指引}
%----------------------------------------------------------

\subsection{STEMI 診斷標準}

根據 2023 ESC Guidelines for Acute Coronary Syndromes：

\begin{longtable}{p{4cm}p{10cm}}
\toprule
\textbf{ST Elevation 標準} & \textbf{定義} \\
\midrule
\endhead
V2-V3 導程 & 男性 $\geq$ 40 歲：$\geq$ 2 mm\newline男性 $<$ 40 歲：$\geq$ 2.5 mm\newline女性：$\geq$ 1.5 mm \\
\midrule
其他導程 & $\geq$ 1 mm（至少連續兩個相鄰導程） \\
\midrule
Posterior MI & V1-V3 ST depression $\geq$ 0.5 mm \\
\bottomrule
\end{longtable}

\subsection{Door-to-Balloon Time 目標}

\begin{importantbox}
\textbf{時間就是心肌！(Time is Muscle)}

\begin{itemize}[noitemsep]
\item \textbf{Door-to-ECG time：} $\leq$ \textbf{10 分鐘}
\item \textbf{Door-to-Balloon time：} $\leq$ \textbf{90 分鐘}（若具備 primary PCI 能力）
\item \textbf{Door-to-Needle time：} $\leq$ \textbf{30 分鐘}（若選擇 fibrinolysis）
\end{itemize}
\end{importantbox}

\subsection{本案例時間分析}

\begin{longtable}{p{5cm}p{9cm}}
\toprule
\textbf{指標} & \textbf{本案例表現} \\
\midrule
\endhead
Door-to-ECG & 約 5 分鐘 \textcolor{green!50!black}{\checkmark} 符合標準 \\
\midrule
ECG 判讀至心導管啟動 & \textcolor{red}{\textbf{未啟動}} $\times$ \\
\midrule
總延誤時間 & $>$ 4 小時（至第四份 ECG 仍未處理） \\
\bottomrule
\end{longtable}

\newpage
%----------------------------------------------------------
\section{改善建議}
%----------------------------------------------------------

\subsection{立即可執行措施}

\begin{enumerate}[leftmargin=*]
\item \textbf{強制閱讀電腦警示}
\begin{itemize}[noitemsep]
\item ECG 出現 ``suspected AMI'' 或 ``STEMI'' 時，系統強制跳出視窗
\item 醫師必須點選確認並選擇處置方向
\end{itemize}

\item \textbf{建立 STEMI 快速通道}
\begin{itemize}[noitemsep]
\item ECG 判讀為 STEMI → 自動通知心臟科值班醫師
\item 啟動心導管室不需等待主治醫師同意
\end{itemize}

\item \textbf{雙重確認機制}
\begin{itemize}[noitemsep]
\item 胸痛病患的 ECG 必須由兩位醫師確認
\item 至少一位需為急診主治醫師層級
\end{itemize}
\end{enumerate}

\subsection{教育訓練建議}

\begin{learningbox}
\textbf{每位急診醫師必須：}
\begin{enumerate}[noitemsep]
\item 熟悉 STEMI 的 ECG 診斷標準
\item 了解電腦判讀的限制與價值
\item \textbf{絕不忽略 ``suspected AMI'' 警示}
\item 建立「寧可過度處理，不可漏診」的心態
\item 定期參加 ECG 判讀訓練
\end{enumerate}
\end{learningbox}

\newpage
%----------------------------------------------------------
\section{Take Home Messages}
%----------------------------------------------------------

\begin{alertbox}
\textbf{永遠不要忽略心電圖機的警示訊息！}

電腦判讀雖有偽陽性，但當警示為\textbf{「suspected AMI」}時：
\begin{enumerate}[noitemsep]
\item \textbf{必須}親自仔細判讀 ECG 波形
\item \textbf{必須}結合臨床症狀綜合判斷
\item \textbf{有疑慮時}立即會診心臟科
\item \textbf{絕不可}直接忽略或等待追蹤
\end{enumerate}
\end{alertbox}

\vspace{1cm}

\begin{keypoint}
\textbf{STEMI 漏診的法律風險：}
\begin{itemize}[noitemsep]
\item 明顯的診斷延誤可能構成醫療疏失
\item 電腦警示紀錄會成為法律證據
\item 「我沒看到」不是有效的辯護理由
\item 病患預後不良將面臨高額賠償
\end{itemize}
\end{keypoint}

\vspace{1cm}

\begin{learningbox}
\textbf{記住這個案例：}

一位 52 歲男性，典型胸痛症狀，心電圖機\textbf{兩次}警示 ``STE over inferior leads, suspected AMI''，Troponin 3 小時內上升 5 倍...

\textbf{但醫療團隊選擇忽略警示。}

這是可以避免的悲劇。
\end{learningbox}

\newpage
%----------------------------------------------------------
\section{附錄：Inferior STEMI 快速判讀口訣}
%----------------------------------------------------------

\subsection{看 ECG 的 30 秒檢查清單}

\begin{enumerate}[leftmargin=*]
\item \textbf{先看電腦判讀} → 有無 ``AMI'', ``STEMI'', ``ischemia'' 等關鍵字
\item \textbf{看 II, III, aVF} → 有無 ST elevation
\item \textbf{看 I, aVL} → 有無 reciprocal ST depression
\item \textbf{結合臨床} → 胸痛 + ST elevation = 啟動心導管！
\end{enumerate}

\subsection{Inferior MI 的血管分布}

\begin{longtable}{p{4cm}p{10cm}}
\toprule
\textbf{梗塞血管} & \textbf{判斷線索} \\
\midrule
\endhead
RCA (Right Coronary Artery) & ST elevation III $>$ II\newline常合併右心室梗塞 (V4R elevation)\newline較常見 (約 80\%) \\
\midrule
LCx (Left Circumflex) & ST elevation II $>$ III\newline可能合併 lateral wall (I, aVL, V5-V6 elevation)\newline較少見 (約 20\%) \\
\bottomrule
\end{longtable}

\vspace{2cm}

\begin{center}
\Large\bfseries\color{emergency_red}
時間就是心肌\\
Time is Muscle\\
\vspace{0.5cm}
\normalsize\color{black}
每延誤 30 分鐘，死亡率增加 7.5\%
\end{center}

\vfill

\begin{center}
\small
本教材由 謝慕揚醫師 MD, PhD, FESC 製作\\
僅供院內教學使用，請勿外流
\end{center}

\end{document}
